
\section{Field Structure}
\subsection{Operational Consequences of the Compatibility Condition}
Expressing the compatibility condition explicitly yields
\begin{equation}\label{eq_compatibility_explicit}
    \dot{\mathbf{e}}_{t}=\dot{\boldsymbol{\omega}}_{t}\mathbf{e}_{0}\boldsymbol{\omega}_{t}^{\dagger}+\boldsymbol{\omega}_{t}\mathbf{e}_{0}\dot{\boldsymbol{\omega}}_{t}^{\dagger}=0,
\end{equation}
while differentiating the unitarity relation $\boldsymbol{\omega}_{t}^{\dagger}\boldsymbol{\omega}_{t}=1$ gives
\begin{equation}
    \dot{\boldsymbol{\omega}}_{t}^{\dagger}=-\boldsymbol{\omega}_{t}^{\dagger}\dot{\boldsymbol{\omega}}_{t}\boldsymbol{\omega}_{t}^{\dagger},
\end{equation}
so substitution into Eq.~\eqref{eq_compatibility_explicit} leads to
\begin{equation}
    \dot{\mathbf{e}}_{t}=\boldsymbol{\omega}_{t}\left[\boldsymbol{\omega}_{t}^{\dagger}\dot{\boldsymbol{\omega}}_{t},\mathbf{e}_{0}\right]\boldsymbol{\omega}_{t}^{\dagger}=0,
\end{equation}
which, owing to rotor invertibility, requires
\begin{equation}\label{eq_initial_commutator}
    \left[\boldsymbol{\omega}_{t}^{\dagger}\dot{\boldsymbol{\omega}}_{t},\mathbf{e}_{0}\right]=0.
\end{equation}
Within the three-dimensional temporal subspace, the bivectors commuting with $\mathbf{e}_{0}$ are proportional to $\mathbf{i}=\mathbf{I}_{t}\mathbf{e}_{0}$, so the compatible angular velocity lies in the plane orthogonal to the fixed causal direction. Since $\mathbf{I}_{t}$ and $\mathbf{e}_{0}$ are fixed in the common reference frame, Eq.~\eqref{eq_pon_normal_fixed} consequently gives $\dot{\mathbf{i}}=0$.

Furthermore, differentiating $\boldsymbol{\omega}_{t}=\cos\theta^{t}-\mathbf{i}\sin\theta^{t}$ for a unit temporal bivector satisfying $\mathbf{i}^{2}=-1$ and right-multiplying by $\boldsymbol{\omega}_{t}^{\dagger}$ yields
\begin{equation}
\begin{split}
    \dot{\boldsymbol{\omega}}_{t}\boldsymbol{\omega}_{t}^{\dagger}={}-\dot{\theta}^{t}\mathbf{i}-\frac{1}{2}\sin(2\theta^{t})\dot{\mathbf{i}}+\sin^{2}\theta^{t}\mathbf{i}\dot{\mathbf{i}},
\end{split}
\end{equation}
where the established constancy of $\mathbf{i}$ eliminates its derivative terms and leaves
\begin{equation}\label{eq_temporal_rotor_evolution}
    \dot{\boldsymbol{\omega}}_{t}\boldsymbol{\omega}_{t}^{\dagger}=\boldsymbol{\omega}_{t}^{\dagger}\dot{\boldsymbol{\omega}}_{t}=-\dot{\theta}^{t}\mathbf{i}.
\end{equation}
Consolidating Eqs.~\eqref{eq_pon_normal_fixed}, \eqref{eq_initial_commutator}, and \eqref{eq_temporal_rotor_evolution}, the compatibility requirement gives three primary operational relations,
\begin{subequations}\label{eq_spinor_compatibility_relations}
\begin{align}
    \dot{\mathbf{i}}=0\implies\dot{\hat{\vartheta}}^{k}&=0,\label{eq_temporal_generator_constancy}\\
    [\mathbf{i},\mathbf{e}_{0}]&=0,\label{eq_temporal_generator_commutator}\\
    [\boldsymbol{\omega}_{t}^{\dagger},\dot{\boldsymbol{\omega}}_{t}]&=0,\\
    \dot{\boldsymbol{\omega}}_{t}+\mathbf{i}\dot{\theta}^{t}\boldsymbol{\omega}_{t}&=0,\label{eq_temporal_rotor_transport}
\end{align}
\end{subequations}
where the first relation expresses constancy along the worldline through $\dot{\hat{\vartheta}}^{k}=U^{\mu}\partial_{\mu}\hat{\vartheta}^{k}$. In the common reference frame, the fixed temporal basis and its direction cosines provide the reference for the field realizations developed below.

\subsection{General Form of the Field}
The temporal rotation is characterized by the scalar angle $\theta^{t}$, its generator, and its composition law, independently of the specific geometric object being transformed. The vector-level action $\mathbf{e}_{t}=\boldsymbol{\omega}_{t}\mathbf{e}_{0}\boldsymbol{\omega}_{t}^{\dagger}$ with $\boldsymbol{\omega}_{t}=\exp(-\mathbf{i}\theta^{t})$ represents one concrete realization of this algebraic symmetry, whose action on other fields follows through their corresponding representations.

Consider an admissible field of type $F$ carrying a unitary representation of the compatible rotor subgroup. Letting $\boldsymbol{\Psi}_{F0}$ denote a fixed reference state and $\mathbf{i}_{F}$ the corresponding constant anti-Hermitian generator, the field along its temporal-rotation orbit is expressed as
\begin{equation}\label{eq_generalized_psi}
    \boldsymbol{\Psi}_{F}=\boldsymbol{\omega}_{tF}\boldsymbol{\Psi}_{F0},\quad \boldsymbol{\omega}_{tF}=\exp(-\mathbf{i}_{F}\theta^{t}),\quad \mathbf{i}_{F}^{\dagger}=-\mathbf{i}_{F},
\end{equation}
with composition and inverse relations
\begin{equation}
\begin{gathered}
    \boldsymbol{\omega}_{tF}(\theta_{1}^{t})\boldsymbol{\omega}_{tF}(\theta_{2}^{t})=\boldsymbol{\omega}_{tF}(\theta_{1}^{t}+\theta_{2}^{t}),\\
    \boldsymbol{\omega}_{tF}^{\dagger}=\boldsymbol{\omega}_{tF}^{-1}.
\end{gathered}
\end{equation}
Single-valued representations of this rotor subgroup inherit its parameter period and therefore satisfy
\begin{equation}
    \exp(-2\pi\mathbf{i}_{F})=1,
\end{equation}
while the transport relation in Eq.~\eqref{eq_temporal_rotor_transport} becomes
\begin{equation}\label{eq_starting_rotor_evolution}
    \dot{\boldsymbol{\omega}}_{tF}+\mathbf{i}_{F}\dot{\theta}^{t}\boldsymbol{\omega}_{tF}=0.
\end{equation}
The fixed reference element thus sweeps out its rotation orbit. In the real vector realization, the same rotation acts by the adjoint transformation
\begin{equation}
\begin{gathered}
    \boldsymbol{\omega}_{t}^{\mathrm{SO}(3)}\mathbf{v}\coloneqq\boldsymbol{\omega}_{t}\mathbf{v}\boldsymbol{\omega}_{t}^{\dagger},\\
    \mathbf{i}_{\mathrm{SO}(3)}\mathbf{v}\coloneqq\mathbf{i}\mathbf{v}-\mathbf{v}\mathbf{i},
\end{gathered}
\end{equation}
so the mappings
\begin{equation}
\begin{gathered}
    \boldsymbol{\Psi}_{F}\to\mathbf{e}_{t},\quad\boldsymbol{\Psi}_{F0}\to\mathbf{e}_{0},\\
    \boldsymbol{\omega}_{tF}\to\boldsymbol{\omega}_{t}^{\mathrm{SO}(3)},\quad\mathbf{i}_{F}\to\mathbf{i}_{\mathrm{SO}(3)}
\end{gathered}
\end{equation}
give $\mathbf{i}_{\mathrm{SO}(3)}\mathbf{e}_{0}=[\mathbf{i},\mathbf{e}_{0}]=0$. Consequently, this real rotation fixes the causal direction while rotating transverse temporal vectors through an angle $2\theta^{t}$ within the plane orthogonal to $\mathbf{e}_{0}$.

Within the common reference frame, the chosen generator satisfies
\begin{equation}
    \dot{\mathbf{i}}_{F}=0,
\end{equation}
while the field realization may admit a larger algebra $\mathfrak{g}_{F}$ containing additional generators. For nonzero $\mathbf{i}_{F}$, choosing a positive-definite inner product in which this generator has unit norm and an orthonormal basis $\{\mathbf{i}_{\beta}\}$ gives
\begin{equation}\label{eq_generator_direction}
\begin{gathered}
    \mathbf{i}_{F}=\hat{\vartheta}^{\beta}\mathbf{i}_{\beta},\quad\sum_{\beta=1}^{\dim\mathfrak{g}_{F}}(\hat{\vartheta}^{\beta})^{2}=1,\\
    \beta\in\{1,\ldots,\dim\mathfrak{g}_{F}\},
\end{gathered}
\end{equation}
where $\hat{\vartheta}^{\beta}$ specifies the generator direction within $\mathfrak{g}_{F}$ and $\theta^{t}$ retains the rotor-angle normalization inherited above. In the temporal vector realization, the three temporal bivectors are dual to their normal directions, so their coefficients correspond to the temporal direction cosines.

\subsubsection{Field Representation and Spin}
The Lie bracket closure
\begin{equation}
    [\mathbf{i}_{\beta},\mathbf{i}_{\gamma}]=f_{\beta\gamma}{}^{\delta}\mathbf{i}_{\delta}
\end{equation}
defines the structure of $\mathfrak{g}_{F}$. Selecting its associated group $G_{F}$ and a unitary representation
\begin{equation}
    \rho_{F}:G_{F}\rightarrow U(V_{F})
\end{equation}
determines the matrix generator
\begin{equation}
    \mathfrak{i}_{F}\coloneqq d\rho_{F}(\mathbf{i}_{F})\in d\rho_{F}(\mathfrak{g}_{F})\subseteq\mathfrak{u}(F),
\end{equation}
and hence the matrix rotor and field state,
\begin{equation}
\begin{gathered}
    w_{F}\coloneqq\rho_{F}(\boldsymbol{\omega}_{tF})=\exp(-\mathfrak{i}_{F}\theta^{t}),\\
    \Psi_{F}=w_{F}\Psi_{F0},
\end{gathered}
\end{equation}
where $V_{F}\cong\mathbb{C}^{F}$ and $\Psi_{F0},\Psi_{F}$ are the corresponding column states. Unitarity then gives $\Psi_{F}^{\dagger}\Psi_{F}=\Psi_{F0}^{\dagger}\Psi_{F0}$, with $\theta^{t}$ describing evolution along the chosen orbit.

Representative field spaces and their admitted algebras include
\begin{equation}
    (V_{F},\mathfrak{g}_{F})\in\bigl\{(\mathbb{C}^{1},\mathfrak{u}(1)),(\mathbb{C}^{2},\mathfrak{su}(2)),(\mathbb{C}^{3},\mathfrak{su}(3))\bigr\},
\end{equation}
where $F=\dim_{\mathbb{C}}V_{F}$ counts the complex field components. The scalar, spinor, and vector realizations are organized by this component count, expressed through the spin index $j$ as
\begin{equation}
    F=2j+1.
\end{equation}
Thus, the one-component scalar field carries the $U(1)$ phase action, the two-component spin-$\frac{1}{2}$ field carries the fundamental action of $SU(2)$, and the three-component spin-$1$ field carries the fundamental action of $SU(3)$. The phase generator $\mathfrak{i}_{F}$ describes the selected temporal rotation, while any additional admitted field transformations belong to the broader algebra.

For the three-component realization, the real rotation subgroup follows by restricting $SU(3)$ to transformations preserving $\mathbb{R}^{3}\subset\mathbb{C}^{3}$, giving
\begin{equation}
    SU(3)\cap GL(3,\mathbb{R})=SO(3),
\end{equation}
or, at the Lie-algebra level,
\begin{equation}
    \mathfrak{su}(3)\cap\mathfrak{gl}(3,\mathbb{R})=\mathfrak{so}(3).
\end{equation}
Accordingly, a real initial vector remains real under the corresponding generator,
\begin{equation}
    \mathfrak{i}_{F}\in\mathfrak{so}(3),\quad\Psi_{F0}\in\mathbb{R}^{3}\implies\Psi_{F}\in\mathbb{R}^{3},
\end{equation}
retaining three real components with the natural vector rotation action. The full $SU(3)$ realization acts on the complex three-component field, while its real $SO(3)$ subgroup supplies this vector realization. Spacetime dependence remains parameterized by $X^{\mu}$, with proper-time evolution evaluated along $X^{\mu}(\tau)$.

\subsubsection{Gauge Covariance and Non-Abelian Fields}\label{ch_gauge_covariance}
Because the phase origin of $\theta^{t}$ is arbitrary, the Principle of Relativity demands that field equations exhibit covariance under local shifts of this origin. Such a transformation acts on the temporal rotor according to
\begin{equation}
\begin{gathered}
    \boldsymbol{\omega}_{tF}'=\exp[-\mathbf{i}_{F}(\theta^{t}+\delta\theta^{t})]=\mathbf{g}\boldsymbol{\omega}_{tF},\\
    \mathbf{g}=\exp(-\mathbf{i}_{F}\delta\theta^{t}),
\end{gathered}
\end{equation}
where $\delta\theta^{t}(X)$ represents the local angular displacement, with $\delta\theta^{t}(\tau)=\delta\theta^{t}(X(\tau))$ along a trajectory.

Expanding $\mathbf{i}_{F}=\hat{\vartheta}^{\beta}\mathbf{i}_{\beta}$ in the admitted generator basis, the local transformation becomes
\begin{equation}
\begin{gathered}
    \mathbf{g}=\exp[-\alpha^{\beta}(X)\mathbf{i}_{\beta}],\\
    g_{F}\coloneqq\rho_{F}(\mathbf{g})=\exp[-\alpha^{\beta}(X)T_{\beta}],
\end{gathered}
\end{equation}
with
\begin{equation}
\begin{gathered}
    T_{\beta}\coloneqq d\rho_{F}(\mathbf{i}_{\beta}),\\
    \alpha^{\beta}(X)=\delta\theta^{t}(X)\hat{\vartheta}^{\beta}\in\mathbb{R}.
\end{gathered}
\end{equation}
These proportional parameters describe the selected phase orbit, while general reference changes within the admitted group use independent $\alpha^{\beta}(X)$. Under this action, the abstract field transforms in tandem with the temporal rotor,
\begin{equation}
\begin{gathered}
    \boldsymbol{\Psi}_{F}\mapsto\mathbf{g}\boldsymbol{\Psi}_{F},\quad w_{F}\mapsto g_{F}w_{F},\\
    \Psi_{F}\mapsto g_{F}\Psi_{F},
\end{gathered}
\end{equation}
so the same covariance applies to every subsequent equation containing derivatives of the temporal rotor. The Principle of Relativity therefore leads to the corresponding gauge-covariant derivative and, through its standard construction, the Abelian or non-Abelian potentials associated with the admitted $U(1)$, $SU(2)$, and $SU(3)$ field structures.
